\part{Linear Elliptic Equations}

\section{Basic definitions and examples}
\subsection{Elliptic operators and principal symbols}

\begin{definition}
Let $P$ be a partial differential operator. The \emph{symbol} of $P$ is obtained by replacing each derivative $\partial_j$ by $i\xi_j$.

Suppose now that $P$ is a linear partial differential operator acting on vector-valued functions $u=(u_J)_{J=1}^N:U\to\R^N$, with leading part of the form
\[
(Pu)^I = \sum_{J,\,|\alpha|=K} A_{J,\alpha}^I(x)\partial^\alpha u_J + \text{(lower order terms)}.
\]
Here $K$ is called the \emph{order} of $P$.
\end{definition}

\begin{definition}
The \emph{principal symbol} of a $K$-order differential operator $P$ is
\[
\sigma_{\mathrm{prin}}(P)(x,\xi)
= i^K\sum_{|\alpha|=K} A^I_{J,\alpha}(x)\xi^\alpha.
\]
We say that $P$ is \emph{elliptic} if $\sigma_{\mathrm{prin}}(P)(x,\xi)$ is invertible for every $x\in U$ and every $\xi\neq 0$.
\end{definition}

\begin{definition}[Uniform ellipticity]
A second-order operator with symmetric coefficient matrix $a=[a_{ij}]$ is \emph{uniformly elliptic} if there exists $\lambda>0$ such that
\[
a_{ij}(x)\xi_i\xi_j \geq \lambda |
\xi|^2
\qquad \text{for all }x\in U,\ \xi\in\R^d.
\]
Equivalently, every eigenvalue of $a(x)$ is bounded below by $\lambda$ uniformly in $x$.
\end{definition}

\begin{definition}[Divergence form operator]
Let $U\subset\R^d$ be open. A second-order operator is said to be in \emph{divergence form} if
\[
Pu = -\partial_i(a^{ij}\partial_j u) + \partial_i(b^i u) + cu,
\]
where $a^{ij}$, $b^i$, and $c$ are coefficient functions.

If the coefficients are smooth enough, one may expand the operator into non-divergence form:
\[
Pu = -a^{ij}\partial_i\partial_j u + \bigl(b^i-\partial_j a^{ij}\bigr)\partial_i u + (\partial_i b^i + c)u.
\]

\end{definition}

\subsection{An $H^1$ a priori estimate}

The basic tool in existence and uniqueness theory is an a priori estimate.

\begin{theorem}[A priori estimate]
Suppose $u\in H_0^1(U)$ solves the Dirichlet problem
\[
\begin{cases}
Pu = f & \text{in }U,\\
u = 0 & \text{on }\partial U,
\end{cases}
\]
where
\[
Pu = -\partial_i(a^{ij}\partial_j u) + \partial_i(b^i u) + cu.
\]
Assume that $P$ is uniformly elliptic, that $a^{ij},b^i,c$ are smooth, and that $a,b,c\in L^\infty(U)$ with
\[
\norm{b}_{L^\infty(U)} + \norm{c}_{L^\infty(U)} \leq A.
\]
Then there exist constants $C>0$ and $\gamma\geq 0$ such that
\[
\norm{u}_{H^1(U)} \leq C\norm{f}_{H^{-1}(U)} + \gamma \norm{u}_{L^2(U)}.
\]
\end{theorem}
\begin{proof}
    For $v \in C^{\infty}_c(U)$, we have
    \[
    \begin{aligned}
    \int_U (Pv)v &= \int_U (-\partial_i(a^{ij}\partial_j v) + \partial_i(b^i v) + cv)v \\
    &= \int_U a^{ij}(\partial_j v)(\partial_i v) - b^i v(\partial_i v) + cv^2 \\
    &\geq \int_U \lambda|D v|^2 - vb\cdot Dv + cv^2
    \end{aligned}
    \]
    by the uniform ellipticity of $P$. Notice by Cauchy inequality, we have
    \[
    \left\lvert \int_U vb\cdot Dv\right\rvert \leq \int_U |b||v||Dv| \leq A\lVert v\rVert_{L^2(U)}\lVert Dv\rVert_{L^2(U)}
    \]
    and
    \[\int_U cv^2 \leq A\lVert v\rVert_{L^2(U)}^2\]
    so we have
    \[
    \begin{aligned}
    \left|\int_U (Pv)v\right| + A\lVert v\rVert_{L^2(U)}\lVert Dv\rVert_{L^2(U)} + A\lVert v\rVert_{L^2(U)} &\geq \left|\int_U (Pv)v\right| + \int_U vb\cdot Dv - cv^2\\
    &\geq \lambda \lVert Dv\rVert_{L^2(U)}^2.
    \end{aligned}
    \]
    Then since $a \in L^{\infty}(U)$, we may let $v$ to approach $u$ in $H^1_0(U)$ and we have
    \[
    ||f||_{H^{-1}(U)}||u||_{H^{1}(U)} + A\lVert u\rVert_{L^2(U)}\lVert Du\rVert_{L^2(U)} + A\lVert u\rVert_{L^2(U)}^2 \geq \lambda \lVert Du\rVert_{L^2(U)}^2  
    \]
    Then use a furthermore estimation, we have
    \[
    ||f||_{H^{-1}(U)}||u||_{H^{1}(U)} + \dfrac A2(t\lVert u\rVert_{L^2(U)}^2 + t^{-1}\lVert Du\rVert_{L^2(U)}^2) + A\lVert u\rVert_{L^2(U)}^2 \geq \lambda \lVert Du\rVert_{L^2(U)}^2 
    \]
    for some $t > 0$ and hence
    \[
     ||f||_{H^{-1}(U)}||u||_{H^{1}(U)} + A(t/2 + 1)\lVert u\rVert_{L^2(U)}^2 \geq (\lambda - A/2t) \lVert Du\rVert_{L^2(U)}^2.
    \]
    Now we have
    \[
    (\lambda - A/2t) \lVert u\rVert_{H^1(U)}^2 \leq ||f||_{H^{-1}(U)}||u||_{H^{1}(U)} + [A(t/2 + 1)+(\lambda - A/2t)]\lVert u\rVert_{L^2(U)}^2,
    \]
    and denote $C = (\lambda - A/2t)^{-1 } > \lambda^{-1}$ and we have
    \[
    \lVert u\rVert_{H^1(U)} \leq C||f||_{H^{-1}(U)} + \gamma\lVert u\rVert_{L^2(U)}
    \]
    where $\gamma :=(\lambda-C^{-1})^{-1}A^2/4 + A + 1 $.
\end{proof}

\subsection{A functional-analytic solvability lemma}

The abstract principle underlying the solvability theory is the following.

\begin{lemma}
Let $X$ and $Y$ be Banach spaces, and let $P:X\to Y$ be bounded and linear.
\begin{enumerate}
    \item If
    \[
    \norm{u}_X \leq C\norm{Pu}_Y,
    \]
    then $\ker P=\{0\}$, and for every $g\in X^*$ there exists $v\in Y^*$ such that
    \[
    P^*v=g,
    \qquad
    \norm{v}_{Y^*}\leq C\norm{g}_{X^*}.
    \]
    \item If
    \[
    \norm{v}_{Y^*} \leq C'\norm{P^*v}_{X^*},
    \]
    then $\ker P^*=\{0\}$, and for every $f\in Y$ there exists $u\in X$ such that
    \[
    Pu=f,
    \qquad
    \norm{u}_X\leq C'\norm{f}_Y.
    \]
\end{enumerate}
\end{lemma}

\begin{lemma}
If $b=0$ and $c=0$, so that
\[
Pu=-\partial_j(a_{jk}\partial_k u),
\]
then the energy estimate holds with $\gamma=0$, namely
\[
\norm{u}_{H^1(U)}\leq C\norm{f}_{H^{-1}(U)}
\]
for solutions $u\in H_0^1(U)$ of $Pu=f$.
\end{lemma}

\begin{theorem}
For every $f\in H^{-1}(U)$, there exists a unique $u\in H_0^1(U)$ such that
\[
-\partial_j(a_{jk}\partial_k u)=f
\qquad \text{in }U.
\]
Thus the homogeneous Dirichlet problem is uniquely solvable in the pure divergence-form case.
\end{theorem}
\begin{proof}
    Let $X = H^1_0(U), Y = H^{-1}(U)$ and notice that they are dual spaces. Applying the above lemma is straight-forward.
\end{proof}

\subsection{Fredholm alternative for general second-order elliptic operators}


\begin{lemma}[Fredholm alternative for compact perturbations]\ \par
(a) For $K:X\to Y$, $K$ is compact iff $K^*$ is compact.\par
(b) Let $K:X\to X$ be compact and set $T=I+K$.
\begin{enumerate}
    \item $\ker(I+K)$ is finite dimensional.
    \item There exists $n_0\geq 1$ such that
    \[
    \ker(I+K)^n = \ker(I+K)^{n_0}
    \qquad \text{for all }n\geq n_0.
    \]
    \item $\ran(I+K)$ is closed and satisfies
    \[
    \ran(I+K) = \bigl(\ker(I+K^*)\bigr)^\perp.
    \]
    \item The dimensions of the kernels agree:
    \[
    \dim \ker(I+K)=\dim \ker(I+K^*).
    \]
\end{enumerate}
\end{lemma}
\begin{theorem}[Fredholm alternative for elliptic Dirichlet problems]
Let $P$ be a second-order scalar elliptic operator on a bounded domain $U$ with $C^1$ boundary. Then exactly one of the following alternatives holds:
\begin{enumerate}
    \item For every $f\in H^{-1}(U)$, there exists a unique $u\in H_0^1(U)$ such that
    \[
    Pu=f,
    \]
    and moreover
    \[
    \norm{u}_{H^1(U)}\leq C\norm{f}_{H^{-1}(U)}.
    \]
    \item There exists a nonzero $u\in H_0^1(U)$ such that
    \[
    Pu=0.
    \]
\end{enumerate}
If the second alternative holds, then both $\ker P$ and $\ker P^*$ are finite dimensional, and the equation
\[
Pu=f
\]
is solvable if and only if
\[
\langle f,v\rangle = 0
\qquad \text{for all }v\in \ker P^*.
\]
\end{theorem}
\begin{proof}
    Let
    \[
    Pu = -\partial_j(a^{jk}\partial_k u)+b^j\partial_j u + cu
    \]
    and we will have the energy estimate
    \[
    \lVert u\rVert_{H^1_0(U)} \leq C\lVert Pu\rVert_{H^{-1}(U)} + \gamma\lVert u\rVert_{L^2(U)}.
    \]
    We consider operator $(P+tI)$, then consider
    \[
    \int_U (P+tI)u^2 \geq \lambda\int_U|Du|^2 + \int_U (b\cdot Du)u + (c+t)u^2 \geq \lambda' \left(\int_U |Du|^2 + u^2\right)
    \]
    when $t$ is large enough for some $\lambda'$, so we have
    \[
    \lVert (P+tI)u\rVert_{H^{-1}(U)}\lVert u\rVert_{H^1_0(U)} \geq \lambda'\lVert u\rVert_{H^1_0(U)}^2
    \]
    and hence $(P+tI)^{-1}$ is well-defined on $H^{-1}(U)$ for $t$ large enough. Then we know
    \[
    u - t(P+tI)^{-1}u = (P+tI)^{-1}f
    \]
    and define the operator
    \[
    T: L^2(U) \xrightarrow{\text{naturally}}H^{-1}(U) \xrightarrow{(P+tI)^{-1}} H^1_0(U) \xrightarrow{\text{compact embedding by Rellich-Kondrachov}} L^2(U)
    \]
    when $n \geq 3$. Then $T$ is compact and hence $\ker (I + t'T) < \infty$. Let $t' = -t$ and then we know if $Pu = 0$, then $u \in \ker(I+t'T)$ and hence $\ker P$ is finite. By the way, we consider the lemma above and it gives $\ker(I+K^*) $
\end{proof}

\section{$L^2$-based elliptic regularity}
\subsection{Interior $H^2$ regularity for general elliptic operators}
\begin{theorem}[$H^2$ elliptic regularity]
Let $u\in H^1(U)$ be a weak solution of
\[
Pu=f
\qquad \text{in }U,
\]
with $f\in L^2(U)$. Assume that $P$ is uniformly elliptic, that $a,b,c\in L^\infty(U)$, and that the coefficients have enough regularity to justify commuting with derivatives. Then for every $V\Subset U$,
\[
u\in H^2(V)
\qquad \text{and} \qquad
\norm{u}_{H^2(V)} \leq C\bigl(\norm{f}_{L^2(U)} + \norm{u}_{L^2(U)}\bigr).
\]
\end{theorem}
\begin{proof}
    
\end{proof}

\begin{definition}
Fix $k\in\{1,\dots,d\}$ and $h\in\R\setminus\{0\}$. The \emph{difference quotient} in the $e_k$ direction is defined by
\[
D_k^h v(x) = \frac{v(x+he_k)-v(x)}{h}.
\]
As $h\to 0$, this converges formally to $\partial_k v$.
\end{definition}

\begin{lemma}
Let $V\Subset U$.
\begin{enumerate}
    \item If $u\in W^{1,p}(U)$, then for $|h|\ll 1$,
    \[
    \norm{D_k^h u}_{L^p(V)} \leq C\norm{\partial_k u}_{L^p(U)}.
    \]
    \item Conversely, if $u\in L^p(U)$ and the difference quotients satisfy a uniform bound
    \[
    \norm{D_k^h u}_{L^p(V)} \leq A
    \qquad \text{for all sufficiently small }h,
    \]
    then $\partial_k u\in L^p(V)$ and
    \[
    \norm{\partial_k u}_{L^p(V)} \leq A.
    \]
\end{enumerate}
\end{lemma}

These facts allow one to repeat the energy estimate at the level of difference quotients and then pass to the limit.

\subsection{Boundary $H^2$ regularity}

Interior estimates alone do not control second derivatives up to the boundary. For boundary regularity one must combine the elliptic equation with the boundary condition and the geometry of $\partial U$.

\begin{theorem}[Boundary $H^2$ regularity]
Let
\[
Lu = -\partial_i(a_{ij}\partial_j u) + b_i\partial_i u + cu = f
\qquad \text{in }U,
\]
where $P$ is uniformly elliptic and the coefficients satisfy the same assumptions as in the interior $H^2$ theorem. Assume in addition that
\begin{enumerate}
    \item $u\in H_0^1(U)$,
    \item $\partial U$ is of class $C^2$.
\end{enumerate}
Then
\[
u\in H^2(U)
\qquad \text{and} \qquad
\norm{u}_{H^2(U)} \leq C\bigl(\norm{f}_{L^2(U)} + \norm{u}_{L^2(U)}\bigr).
\]
\end{theorem}

\begin{remark}
Near a boundary point, one first straightens the boundary by a $C^2$ change of coordinates, reducing the problem locally to a half-ball. Tangential second derivatives are then estimated by the interior argument, using that the Dirichlet condition forces tangential derivatives to vanish on the flattened boundary. The remaining normal second derivative is recovered from the equation itself, together with uniform ellipticity, which guarantees control of the coefficient of $\partial_d^2 u$.
\end{remark}
